\documentclass[11pt]{article}
\usepackage[utf8]{inputenc}
\usepackage{amsmath,amsthm,amssymb}
\usepackage{geometry}
\geometry{margin=2.5cm}
\usepackage[hidelinks]{hyperref}
\usepackage{booktabs}

\newtheorem{theorem}{Theorem}
\newtheorem{lemma}[theorem]{Lemma}
\newtheorem{corollary}[theorem]{Corollary}
\newtheorem{proposition}[theorem]{Proposition}
\theoremstyle{definition}
\newtheorem{definition}[theorem]{Definition}
\theoremstyle{remark}
\newtheorem{remark}[theorem]{Remark}

\newcommand{\floor}[1]{\left\lfloor #1 \right\rfloor}
\newcommand{\dnorm}[1]{\left\| #1 \right\|}

\title{One-sided regularity of the Beatty ballot constant:\\
left jumps, right flatness, and a closed form for the jump}
\author{Jan Popelka\thanks{Email: popojan@protonmail.com}}
\date{June 2026}

\begin{document}
\maketitle

\begin{abstract}
Let $C(\alpha)$ denote the asymptotic constant of monotone lattice paths
under the Beatty staircase of slope $\alpha > 1$, as studied in the
companion survey~\cite{bridge}.
Working at the level of the underlying Sturmian ruin probabilities, we
prove that at every rational $p/q$ the left limit $C^-(p/q)$ exists and
equals the constant of the \emph{same} periodic walk with a sharpened
absorbing barrier (level $0$ at phase $0$); this yields an algebraic
closed form for the jump $J(p/q) = C(p/q) - C^-(p/q)$.
We further prove that $C$ is flat to infinite order from the right at
every rational, with an explicit exponential bound
$C(p/q + \varepsilon) - C(p/q) \leq A\,e^{-\kappa/\varepsilon}$, and
that $C$ is continuous at every irrational, with a stretched-exponential
modulus governed by the Diophantine type.
For integer slopes the closed form collapses to
$1 - 2\,C^-(k) = \rho_k^{\,k}$, where $\rho_k$ is the ruin probability
of the $\{+k, -1\}$ walk.
\end{abstract}

\section{Setup}

We use the framework of~\cite{bridge}.
Fix a slope $\alpha > 1$.
An infinite sequence of independent fair coin flips generates a path of
right-steps ($R$) and up-steps ($U$); after $t$ flips containing
$\tau_t$ right-steps and $u_t$ up-steps, the \emph{gap process} is
\begin{equation}\label{eq:gap}
s_t \;=\; \floor{\alpha(\tau_t + 1)} - u_t .
\end{equation}
The walk \emph{survives} if $s_t \geq 0$ for all $t$; write
$\sigma(\alpha)$ for the survival probability.
For rational $\alpha = p/q$ the rise sequence
$r_j = \floor{\alpha(j+1)} - \floor{\alpha j}$ is $q$-periodic, the ruin
probability $\rho(s, j)$ from height $s$ in phase $j = x \bmod q$
satisfies the periodic recursion of~\cite{bridge}, and
$\sigma(p/q) = 1 - \rho(s_0, j_0)$ with $s_0 = \floor{p/q}$, $j_0 = 1$.
Throughout we use the identification
\begin{equation}\label{eq:ident}
C(\alpha) = \tfrac{1}{2}\,\sigma(\alpha)
\end{equation}
of~\cite{bridge} (eq.~(C-rational) there; verified against exact path
counts to 20+ digits).
All results below are proved for $\sigma$; they transfer to $C$
via~\eqref{eq:ident}.

Two elementary observations are used repeatedly.
First, since $\floor{x} > x - 1$,
\begin{equation}\label{eq:linbound}
s_t \;>\; \alpha\,\tau_t - u_t + (\alpha - 1) \;=:\; Y_t + (\alpha - 1),
\end{equation}
where $Y_t = \sum_{i \leq t} \zeta_i$ with i.i.d.\ increments
$\zeta_i = +\alpha$ or $-1$, each with probability $\tfrac12$.
Second, monotonicity: if $\alpha' \leq \alpha''$ then
$\floor{\alpha' m} \leq \floor{\alpha'' m}$ for all $m$, so under the
common coin-flip coupling the gap processes satisfy
$s'_t \leq s''_t$ pointwise and $\sigma(\alpha') \leq \sigma(\alpha'')$.

\section{Two lemmas}

\begin{lemma}[Column comparison]\label{lem:columns}
Let $\alpha = p/q$ in lowest terms and $\varepsilon > 0$.
\begin{enumerate}
\item[(a)] For every $m \geq 1$:
  $\floor{(\alpha - \varepsilon)m} \leq \floor{\alpha m} - [\,q \mid m\,]$,
  with equality whenever $\varepsilon m < 1/q$.
\item[(b)] For every $m$ with $\varepsilon m < 1/q$:
  $\floor{(\alpha + \varepsilon)m} = \floor{\alpha m}$.
\item[(c)] For irrational $\alpha$ and any $x$ with $|x - \alpha| =
  \varepsilon$: $\floor{x m} = \floor{\alpha m}$ for all
  $m < m_*(\varepsilon) := \min\{m \geq 1 : \dnorm{\alpha m} \leq
  \varepsilon m\}$, where $\dnorm{\cdot}$ is the distance to the nearest
  integer; $m_*(\varepsilon) \to \infty$ as $\varepsilon \to 0$.
\end{enumerate}
\end{lemma}

\begin{proof}
(a) The fractional part of $\alpha m$ lies in
$\{0, 1/q, \ldots, (q-1)/q\}$ and vanishes exactly when $q \mid m$.
Subtracting $\varepsilon m > 0$ from an integer drops the floor by at
least one; subtracting less than $1/q$ from a number with fractional
part $\geq 1/q$ leaves the floor unchanged.
(b) Adding $\varepsilon m < 1/q$ cannot push the fractional part
($\leq (q-1)/q$) past $1$.
(c) Immediate from the definition; $m_* \to \infty$ because
$\dnorm{\alpha m} > 0$ for each fixed $m$.
\end{proof}

\begin{lemma}[Late-contact bound]\label{lem:chernoff}
Fix $\alpha > 1$ and $\theta \in (0, \theta_0(\alpha))$, where
$\theta_0$ is the positive root of
$\lambda_\theta := \tfrac12(e^{-\theta\alpha} + e^{\theta}) = 1$.
Then for every $c \geq 0$ and $T \geq 1$,
\[
P\bigl(\exists\, t \geq T:\ s_t \leq c\bigr)
\;\leq\; \frac{e^{\theta(c - \alpha + 1)}}{1 - \lambda_\theta}\;
\lambda_\theta^{\,T}.
\]
The same bound holds for any gap process whose allowance differs from
\eqref{eq:gap} by a bounded amount $b$, with $c$ replaced by $c + b$.
\end{lemma}

\begin{proof}
By~\eqref{eq:linbound}, $s_t \leq c$ implies $Y_t \leq c - \alpha + 1$.
For fixed $t$, Markov's inequality on $e^{-\theta Y_t}$ gives
$P(Y_t \leq c - \alpha + 1) \leq e^{\theta(c - \alpha + 1)}
\lambda_\theta^{\,t}$, since
$E\,e^{-\theta\zeta} = \lambda_\theta$.
Note $\lambda_\theta < 1$ for small $\theta > 0$ because
$\frac{d}{d\theta}\lambda_\theta|_{0} = \tfrac12(1 - \alpha) < 0$.
Sum the geometric series over $t \geq T$.
\end{proof}

Since every column (right-step) consumes at least one coin flip,
``after column $M$'' implies $t \geq M$, and
Lemma~\ref{lem:chernoff} applies with $T = M$.

\section{Left limit: the sharpened-barrier system}

\begin{definition}\label{def:sharp}
For $\alpha = p/q$, the \emph{sharpened system} is the periodic walk
with gap process
$\tilde s_t = \floor{\alpha(\tau_t+1)} - [\,q \mid (\tau_t+1)\,] - u_t$,
i.e.\ the original walk with the absorbing barrier raised by one unit
in the columns $q \mid m$ (phase $0$).
Write $\tilde\sigma(p/q)$ for its survival probability and
$\tilde\rho(s, j)$ for its ruin probabilities.
\end{definition}

\begin{proposition}\label{prop:sharp-solve}
$\tilde\rho$ satisfies the same interior recursion as $\rho$, with the
same sub-unit master-equation roots $t_i$ and phase amplitudes
$A_j^{(i)}$ of~\cite{bridge}; the boundary system changes in exactly
one row: the phase-$0$ condition $\sum_i c_i A_0^{(i)}/t_i = 1$ is
replaced by $\sum_i c_i A_0^{(i)} = 1$ (absorption at level $0$
instead of $-1$).
Consequently $\tilde\sigma(p/q) = 1 - \tilde\rho(s_0, j_0)$ is
algebraic.
For $q = 1$ the system collapses to $\tilde\rho(s) = \rho^{\,s}$.
\end{proposition}

\begin{proof}
The interior recursion (conditioning on one coin flip) is untouched by
the barrier modification, so the exponential ansatz and the master
equation are unchanged.
The absorbing states are $s = -1$ at phases $j \neq 0$ and $s = 0$ at
phase $0$, giving the stated $q$ boundary rows.
For $q = 1$: the single condition $\tilde\rho(0) = 1$ with
$\tilde\rho(s) = c\,t^s$, $t = \rho$, forces $c = 1$.
\end{proof}

\begin{theorem}[Left limit]\label{thm:left}
For every rational $p/q$ with $p > q \geq 1$, $\gcd(p,q) = 1$:
\[
\lim_{\varepsilon \to 0^+} \sigma(p/q - \varepsilon)
= \tilde\sigma(p/q),
\qquad\text{with}\qquad
0 \leq \tilde\sigma - \sigma(p/q - \varepsilon)
\leq A\,\lambda^{\,1/(q\varepsilon)}
\]
for explicit $A, \lambda < 1$ depending only on $p/q$.
Hence, via~\eqref{eq:ident},
$C^-(p/q) = \tfrac12\bigl(1 - \tilde\rho(s_0, j_0)\bigr)$ and the jump
$J(p/q) = \tfrac12\bigl(\tilde\rho(s_0,j_0) - \rho(s_0,j_0)\bigr)$ is
algebraic.
\end{theorem}

\begin{proof}
Couple all walks on the same coin flips.
By Lemma~\ref{lem:columns}(a), $s^{(\varepsilon)}_t \leq \tilde s_t$
pointwise, so $\sigma(p/q - \varepsilon) \leq \tilde\sigma$; the left
limit (which exists by monotonicity of $\sigma$ and is independent of
the approach) is $\leq \tilde\sigma$.

For the reverse inequality, set $M = \floor{1/(q\varepsilon)}$.
By Lemma~\ref{lem:columns}(a) the two gap processes coincide while
$\tau_t < M$.
On the event $\{\tilde s \text{ survives}\} \setminus
\{s^{(\varepsilon)} \text{ survives}\}$, ruin of $s^{(\varepsilon)}$
occurs at some time $t$ with $\tau_t \geq M$, at which moment
$\tilde s_t \leq s^{(\varepsilon)}_t + (\varepsilon \tau_t + 1)
\leq \varepsilon\tau_t + 1$
(the two allowances differ by at most $\varepsilon m + 1$ at column
$m$).
Partition the columns into dyadic blocks $\tau \in [2^i M, 2^{i+1} M)$:
on block $i$ the defect is at most $2^{i+1}\varepsilon M + 1 \leq
2^{i+1}/q + 1$, while $t \geq 2^i M$.
Lemma~\ref{lem:chernoff} bounds the probability of block-$i$ contact by
$B\,e^{\theta(2^{i+1}/q + 1)}\lambda_\theta^{\,2^i M}$.
Choosing $\theta$ with $e^{2\theta/q} \lambda_\theta^{M} < 1$ (possible
for $M \geq M_0(\theta)$), the sum over $i$ is dominated by a constant
multiple of its first term, giving the bound
$A \lambda_\theta^{\,M}$.
\end{proof}

\begin{corollary}[Integer slopes]\label{cor:integer}
$1 - 2\,C^-(k) = \rho_k^{\,k}$, i.e.\
$C^-(k) = (\tau_k - 1)/2$ and
$J(k) = (\tau_k - 1)/(2\tau_k^{\,k+1}) \sim 2^{-(k+2)}$,
where $\tau_k = 1/\rho_k$ is the multinacci constant.
\end{corollary}

\begin{proof}
Proposition~\ref{prop:sharp-solve} with $q = 1$:
$\tilde\rho(s_0) = \rho^{\,k}$ since $s_0 = k$; then
$2 - \tau_k = \tau_k^{-k}$ (master equation) gives the stated forms.
\end{proof}

\section{Right flatness}

\begin{theorem}[Right flatness]\label{thm:right}
For every rational $p/q$ and $0 < \varepsilon < 1$:
\[
0 \;\leq\; \sigma(p/q + \varepsilon) - \sigma(p/q)
\;\leq\; A\,\lambda^{\,1/(q\varepsilon)} .
\]
Consequently all one-sided right derivatives of $C$ at $p/q$ vanish:
$C$ is flat to infinite order from the right.
\end{theorem}

\begin{proof}
Monotonicity gives the left inequality.
By Lemma~\ref{lem:columns}(b) the gap processes coincide while
$\tau_t < M = \floor{1/(q\varepsilon)}$, so on
$\{s^{(+\varepsilon)} \text{ survives}\} \setminus
\{s \text{ survives}\}$ the unperturbed walk is ruined at a column
$\geq M$ while carrying a bonus of at most $\varepsilon\tau + 1$;
the dyadic-block argument of Theorem~\ref{thm:left} applies verbatim
to the unperturbed walk's late contact, giving the bound.
Flatness: $A e^{-\kappa/\varepsilon}/\varepsilon^n \to 0$ for every $n$,
with $\kappa = |{\ln \lambda}|/q$.
\end{proof}

\begin{remark}
The same argument applied below the jump gives
$\tilde\sigma - \sigma(p/q - \varepsilon) \leq A\lambda^{1/(q\varepsilon)}$
(Theorem~\ref{thm:left}): the approach to the \emph{left limit} is also
exponentially fast.
Near each rational, $C$ is therefore a two-level step function up to
exponentially small corrections --- consistent with the geometric
convergence (in $K$) of the continued-fraction ladders observed
numerically.
\end{remark}

\section{Continuity at irrationals}

\begin{theorem}\label{thm:irrational}
$\sigma$ (hence $C$) is continuous at every irrational $\alpha$, with
\[
|\sigma(x) - \sigma(\alpha)| \;\leq\; A\,\lambda^{\,m_*(|x - \alpha|)},
\]
$m_*$ as in Lemma~\ref{lem:columns}(c).
If $\alpha$ has Diophantine type $\kappa$ (i.e.\
$\dnorm{\alpha m} \geq c\,m^{1-\kappa}$), then
$m_*(\varepsilon) \geq (c/\varepsilon)^{1/\kappa}$ and the modulus is
stretched-exponential; for bounded type (e.g.\ quadratic irrationals),
$|\sigma(x) - \sigma(\alpha)| \leq A\,\lambda^{\,c/\sqrt{\varepsilon}}$.
\end{theorem}

\begin{proof}
Identical coupling: the words agree up to column $m_*$, the defect/bonus
beyond column $m$ is at most $\varepsilon m + 1$, and the dyadic
late-contact argument gives the bound.
\end{proof}

\section{Numerical verification}

The closed form of Theorem~\ref{thm:left} was tested against one-sided
ladder extrapolations at 25 rationals (all slopes measured in the
companion session); it agrees within each measurement's accuracy,
down to $6 \times 10^{-11}$ at $8/5$ and $4 \times 10^{-9}$ at $5/3$.
Exact example: for $3/2$ (sub-unit roots $0.43691$, $0.71974$, both
real),
\[
J(3/2) = 0.0272922840\ldots,
\qquad
64x^6 - 896x^5 + 4336x^4 - 7688x^3 + 2100x^2 - 528x + 13 = 0 .
\]

The exponential rates of Theorems~\ref{thm:left}--\ref{thm:right} are
visible as geometric convergence of the CF-extension ladders
$x_K \to p/q$ (where $\varepsilon_K = 1/(q\,q_K)$, so
$\lambda^{1/(q\varepsilon)} = \lambda^{q_K}$): measured per-column
contraction $\mu$ versus the per-flip Chernoff bound
$\lambda_{\theta^*}$:

\begin{center}
\begin{tabular}{lcc}
\toprule
slope & measured $\mu$ (above-ladder, $q_K \leq 83$) & bound $\lambda_{\theta^*}$ \\
\midrule
$3/2$ & $0.929$--$0.934$ & $0.9801$ \\
$5/3$ & $0.899$--$0.903$ & $0.9689$ \\
$4/3$ & $0.955$--$0.960$ & $0.9898$ \\
\bottomrule
\end{tabular}
\end{center}

\smallskip\noindent
(The measured rates drift slowly upward with $K$, suggesting a nearby
asymptotic value; all are comfortably below the proven bound.)

The true contraction rate $\mu$ is a spectral quantity of the
time-direction transfer operator; identifying it in closed form is
open.

\begin{remark}[Status]
Theorems~\ref{thm:left}--\ref{thm:irrational} are rigorous statements
about the survival probabilities $\sigma$, $\tilde\sigma$ of the
coin-flip gap processes.
Their translation to the lattice-path constant $C$ rests on the
identification~\eqref{eq:ident} of~\cite{bridge}.
The nonsingularity of the (modified) $q \times q$ boundary system is
assumed as in~\cite{bridge}, where it is verified computationally.
\end{remark}

\begin{thebibliography}{9}
\bibitem{bridge}
J.~Popelka,
\emph{Lattice paths under Beatty staircases: ruin theory and multinacci
constants} (working title),
companion survey, \texttt{ruin-multinacci-bridge}, 2026.
\end{thebibliography}

\end{document}
